This section describes the retesting approach that will be used to
verify the remediation of vulnerabilities identified during the
assessment.

Retesting should be performed after the recommended security fixes
have been implemented. The purpose of retesting is to determine
whether the previously identified vulnerabilities have been
successfully addressed and whether the implemented fixes introduce
any new security issues.

\subsection*{Retesting Objectives}

The main objectives of the retesting activity are:

\begin{itemize}

    \item Verify that the previously confirmed vulnerabilities have
    been remediated.

    \item Reproduce the original test conditions and verify that the
    previously successful proof-of-concept no longer works.

    \item Confirm that the implemented security controls are enforced
    correctly on the server side.

    \item Verify that remediation does not negatively affect the
    intended application functionality.

    \item Identify any vulnerabilities that remain partially or
    completely unresolved.

\end{itemize}

\subsection*{Retesting Methodology}

The retesting process will follow the same general approach used
during the original assessment.

\begin{enumerate}

    \item Review the remediation implemented for each confirmed
    vulnerability.

    \item Reproduce the original request, input, or attack condition
    used during the initial assessment.

    \item Verify whether the previously demonstrated exploitation
    remains possible.

    \item Verify that the relevant security control is now correctly
    enforced.

    \item Test the affected application functionality to confirm that
    the remediation has not introduced unintended behavior.

    \item Record the retesting result and supporting evidence for each
    vulnerability.

\end{enumerate}

\subsection*{Retesting Status}

At the time of preparing this report, remediation had not yet been
implemented against the target application. Retesting was therefore
not performed as part of this engagement and is documented here as
planned future work, to be carried out once the recommended fixes in
Section~\ref{sec:remediation} have been applied.

The following table defines the retesting status that will be used
for the confirmed findings once remediation is complete:

\begin{table}[H]
\centering
\renewcommand{\arraystretch}{1.3}
\begin{tabular}{|l|p{6cm}|l|}
\hline
\textbf{Finding ID} & \textbf{Vulnerability} & \textbf{Status} \\
\hline
BAC &
Broken Access Control -- Privilege Escalation via Registration &
Pending \\
\hline
IDOR &
Insecure Direct Object Reference -- Unauthorized Basket Access &
Pending \\
\hline
AUTH &
Weak Password Recovery &
Pending \\
\hline
SQLI-01 &
SQL Injection -- Authentication Bypass &
Pending \\
\hline
SQLI-02 &
SQL Injection -- Product Search and Data Extraction &
Pending \\
\hline
XSS &
Reflected Cross-Site Scripting -- Search Parameter &
Pending \\
\hline
CSRF &
Cross-Site Request Forgery -- Profile Update &
Pending \\
\hline
OR &
Open Redirect -- Unvalidated Redirect Parameter &
Pending \\
\hline
\end{tabular}
\caption{Retesting status of confirmed vulnerabilities}
\end{table}

\subsection*{Retesting Evidence}

For each remediated vulnerability, evidence should be collected to
demonstrate the result of the retest.

The evidence may include:

\begin{itemize}

    \item HTTP requests and responses captured using Burp Suite.

    \item Screenshots showing the security control preventing the
    previously demonstrated attack.

    \item Updated application behavior after remediation.

    \item Scanner results, where automated verification is
    applicable.

    \item Any remaining error messages or security-related responses.

\end{itemize}

\subsection*{Retesting Result Classification}

Each vulnerability should be assigned one of the following results
after retesting:

\begin{itemize}

    \item \textbf{Remediated} -- The vulnerability could no longer be
    reproduced and the relevant security control was verified.

    \item \textbf{Partially Remediated} -- The original exploitation
    was reduced or restricted, but the vulnerability or its security
    impact remains under certain conditions.

    \item \textbf{Not Remediated} -- The original vulnerability could
    still be reproduced after remediation.

    \item \textbf{Not Retested} -- Retesting could not be performed
    because the required remediation or testing conditions were not
    available.

\end{itemize}

\subsection*{Final Retesting Statement}

Retesting results will be updated after the recommended remediation
activities have been implemented and verified. As this engagement
was conducted against an intentionally vulnerable training
application rather than a production system, remediation and
retesting fall outside the scope of the current assessment and are
recommended as a follow-up activity.

The final retesting record should include the original finding ID,
remediation status, retesting procedure, observed result, supporting
evidence, and final status for each vulnerability.

This approach provides a clear comparison between the original
security assessment and the post-remediation security state of the
application.